\uuid{CYoU}
\titre{Matrice/Inversibilité/Rang}
\auteur{hanani ; mzari}
\organisation{exo7}
\licence{CC BY-NC-SA 4.0}
\datecreate{2020-01-20}
\niveau{L1}
\module{Algèbre}
\chapitre{Matrice}
\sousChapitre{Inverse, méthode de Gauss}
\difficulte{2}
\theme{Inverse d'une matrice, Matrice, Inversibilité, Rang}
\exo7qcmid{lille-2.matrices.026}

\contenu{
\begin{question}
\qclassification{108.05}
\qtags{motcle=Matrice/Inversibilité/Rang}
Soit $ M_n(\Rr)$ l'ensemble des matrices carrées d'ordre $n$ à coefficients réels et $I$ la matrice identité. Soit $A \in M_n(\Rr)$ telle qu'il existe un entier $m\ge 1$ vérifiant $A^m = I$. Quelles sont les assertions vraies ?
\begin{answers}
\good{$A$ est inversible et $A^{-1} = A^{m-1}$}.
\good{Le rang de $A$ est $n$}.
\bad{$A$ n'est pas inversible}.
\good{Si $m=2$, $A$ est inversible et $A^{-1} = A$}.
\end{answers}
\begin{explanations} On a : $A \times A^{m-1} = I$, donc $A$ est inversible et $A^{-1} = A^{m-1}$. Comme $A$ est inversible, le rang de $A$ est $n$. Si $m=2$, alors $A^{-1}=A$.
\end{explanations}
\end{question}
}
